% Source-aligned English for printed p.179.
\SourcePage{179}

\LayerTarget{tohoku:eq:special-cover-H1-is-Z:p2}

\LayerTarget{tohoku:section:3.9}
\noindent\textbf{3.9. Acyclicity criteria by the method of
coverings.}\textsuperscript{8)}

\LayerTarget{tohoku:setup:acyclicity-cover-families}
Let $X$ be a topological space and let $\mathcal S$ be a nonempty collection
of subsets of $X$. For every $A\in\mathcal S$, suppose that a nonempty set
$\mathfrak R(A)$ of coverings $\mathbf R$ of $A$ has been given, where the
members of each covering and their finite intersections belong to
$\mathcal S$. Assume that if $\mathbf R,\mathbf R'\in\mathfrak R(A)$ and
$B\in\mathbf R$, then the trace $\mathbf R'_B$ of $\mathbf R'$ on $B$
belongs to $\mathfrak R(B)$.

\LayerTarget{tohoku:conditions:acyclicity-cover-families}
Assume in addition that one of the following three conditions holds (each of
which allows the Leray--Cartan spectral sequence to be written for every
covering $\mathbf R\in\mathfrak R(A)$, with $A\in\mathcal S$): (i) the
members of $\mathcal S$ are open; (ii) the members of $\mathcal S$ are closed
and the coverings $\mathbf R\in\mathfrak R(A)$ are finite; (iii) the members
of $\mathcal S$ are closed, $X$ is paracompact, and the coverings
$\mathbf R\in\mathfrak R(A)$ are locally finite.

\LayerTarget{tohoku:theorem:3.9.1}
\noindent\textsc{Theorem 3.9.1.} \textit{Under the preceding conditions, let
an abelian sheaf $F$ on $X$ and a natural number $n\geq0$ be given. Suppose
the following conditions hold:}

\LayerTarget{tohoku:condition:3.9.1:A-n}
\begin{quote}\itshape
$A(n)$: $H^i(\mathbf R,F)=0$ for $1\leq i\leq n$ and every
$\mathbf R\in\mathfrak R(A)$, $A\in\mathcal S$.
\end{quote}

\LayerTarget{tohoku:condition:3.9.1:B-n-minus-1}
\begin{quote}\itshape
$B(n-1)$: For every $A\in\mathcal S$ and every $c^i\in H^i(A,F)$, with
$1\leq i\leq n-1$, there is a finite subset $L$ of $\mathfrak R(A)$ such
that, if $\mathbf R^L$ denotes the ``intersection'' covering of the coverings
$\mathbf R\in L$, the restriction of $c^i$ to every member
$B\in\mathbf R^L$ is zero.
\end{quote}

\LayerTarget{tohoku:conclusion:3.9.1}
\noindent\textit{Under these conditions, $H^i(A,F)=\{0\}$ for every
$A\in\mathcal S$ and $1\leq i\leq n-1$. Moreover, if $c^n\in H^n(A,F)$,
then $c^n$ is zero if and only if there is a finite subset $L$ of
$\mathfrak R(A)$ such that the restriction of $c^n$ to every
$B\in\mathbf R^L$ is zero.}

Let us state at once the most interesting corollaries.

\LayerTarget{tohoku:corollary:3.9.1:1}
\noindent\textsc{Corollary 1.} \textit{With the notation of Theorem 3.9.1,
one has $H^i(A,F)=\{0\}$ for every $A\in\mathcal S$ and
$1\leq i\leq n$ if and only if conditions $A(n)$ and $B(n)$ hold.}

\LayerTarget{tohoku:proof:corollary-3.9.1-1:sufficiency}
Indeed, sufficiency follows immediately from the theorem.

\LayerTarget{tohoku:proof:corollary-3.9.1-1:necessity}
Conversely, suppose that $H^i(A,F)=\{0\}$ for $A\in\mathcal S$ and
$1\leq i\leq n$. Then $B(n)$ holds trivially; let us verify $A(n)$. The
Leray spectral sequence for the covering $\mathbf R$ of $A$ (Theorem 3.8.1
and the remark in 3.8) abuts to $H^*(A,F)$ and has initial term
\[
E_2^{p,q}=H^p(\mathbf R,H^q(F)).
\]
This is zero for $1\leq q\leq n$, since
$C(\mathbf R,H^q(F))=0$ for those values of $q$: the finite intersections of
members of $\mathbf R$ belong to $\mathcal S$. It follows in the usual way
that $H^i(A,F)=H^i(\mathbf R,F)$ for $0\leq i\leq n$, and hence
$H^i(\mathbf R,F)=0$ for these $i$.

\LayerTarget{tohoku:corollary:3.9.1:2}
\noindent\textsc{Corollary 2.} \textit{Suppose, with the notation of Theorem
3.9.1, that for every $A\in\mathcal S$ and every open covering of $A$ there
is a finer covering of the form $\mathbf R^L$, where $L$ is a finite subset
of $\mathfrak R(A)$. Then $H^i(A,F)=\{0\}$ for every $A\in\mathcal S$ and
$1\leq i\leq n$ if and only if $H^i(\mathbf R,F)=\{0\}$ for every
$\mathbf R\in\mathfrak R(A)$, with $A\in\mathcal S$, and
$1\leq i\leq n$.}

\LayerTarget{tohoku:proof:corollary-3.9.1-2}
Indeed, condition $B(n)$ holds for every $n>0$ by Lemma 3.8.2; it therefore
suffices to apply Corollary 1.

\TohokuContinuation{180}
